\section{Discussion}
\label{sec:discussion}

\paragraph{The result is not ``DPO beats GRPO''.} It would be easy to read
\S\ref{sec:gains}--\S\ref{sec:heldout} as a verdict on the loss function. It is not. We probed
the update directions of both methods directly by re-weighting each method's own candidate
groups under the \emph{other} method's rule (Appendix~\ref{app:probe}): swapping the rule while
holding the data fixed barely moves the update direction (cosine $\geq0.91$), while holding the
rule fixed across each method's own groups leaves the two as far apart as they were as trained
($0.32$--$0.40$ ceiling-corrected, versus $0.317$ as trained). The two losses extract nearly the
same direction from the same eight completions.

The gap is therefore about \textbf{which states the training data is drawn from}
(\S\ref{sec:setup}), and that suggests a mechanism the outcome scores alone do not:
\textbf{GRPO's training prompts drift along with its own policy.} As the policy becomes more
effusive, the on-policy states it is next trained on are themselves more effusive, and the
update chooses among candidates continued from an already-drifted prefix. PTO's trunk is
re-anchored at every step by an oracle selection, which bounds how far its state distribution
can wander. We offer this as the reading most consistent with our measurements, not a
demonstrated cause: isolating it needs an arm that reranks states while holding the loss fixed,
which we did not run. The drift itself is a \emph{compounding loop, not a strong pull}:
per-iteration selection pressure toward affirmation never exceeds ${\approx}0.10$ while the
policy's generated affirmation rate moves an order of magnitude further
(Appendix~\ref{app:probe}) --- which is why selection-time contrasts understate the drift,
why single-iteration diagnostics miss it, and why retention, which compares policies rather
than candidates, sees it first.

\paragraph{Efficiency changes the comparison, not the conclusion.} An iteration is not a unit
of cost: PTO's ten iterations cost $8.1$ GPU-hours to GRPO's $27.9$ ($3.4\times$ less --- its
tree build runs once per iteration; GRPO recomputes its reward inside the training loop), so on
the compute axis PTO dominates outright. But at \emph{matched} spend (${\approx}8$ GPU-h, GRPO
reaching only iteration 3), PTO@10 wins the reward ($+0.27$, $\dz{}=0.53$) while being markedly
\emph{worse} on MI-inconsistency ($+0.26$, $\dz{}=0.90$; both replicate held-out): the hack
accumulates with optimization depth, so more reward per GPU-hour buys more reward-hacking per
GPU-hour with it.

\paragraph{Practical recommendation.} Whenever an LLM judge supplies the reward, score a
held-out grader from a different family on states you have already generated and report gain
retention per iteration. It costs no human annotation, needs no agreement on absolute level, and
in our runs separated a hacking arm from a non-hacking one four iterations before the reward
curve did. We would treat a retention series falling below ${\approx}0.7$ and trending down as a
stop signal, in the spirit of a validation-loss criterion --- though on one experiment that
threshold is a suggestion, not a calibrated rule. Pair it with a deterministic cross-check of
any behaviour the judge also counts (Figure~\ref{fig:qcross}): disagreement is itself a signal,
and nearly free.

\paragraph{For MI specifically.} The endpoint models are more pleasant and less therapeutic:
they validate more, ask less, and talk longer. A satisfaction-shaped reward selects for a
satisfaction-shaped therapist; if the goal is MI fidelity, the reward has to contain the
prohibition. MI-inconsistency and question rate are cheap to score and were, here, only ever
measured --- putting them in the reward is the natural next experiment, though that also
removes them from the independent checks.
